% Part: incompleteness
% Chapter: arithmetization-syntax
% Section: substitution

\documentclass[../../../include/open-logic-section]{subfiles}

\begin{document}

\olfileid[hi]{inc}{art}{sub}
\olsection{प्रतिस्थापन}

स्मरण करें कि प्रतिस्थापन वह संक्रिया है जिसमें !!a{formula}~$!A$ में
!!a{variable}~$u$ की सभी मुक्त उपस्थितियों को पद~$t$ से बदला जाता है; इसे
$\Subst{!A}{t}{u}$ लिखते हैं। !!{variable}, !!{formula} और पदों की ग्योडेल (G\"odel)
 संख्याओं पर की गई यह संक्रिया आदिम पुनरावर्ती है।

\begin{prop}\ollabel{prop:subst-primrec}
ऐसा आदिम पुनरावर्ती फलन $\fn{Subst}(x, y, z)$ है जिसका गुण है कि
\[
\fn{Subst}(\Gn{!A}, \Gn{t}, \Gn{u}) = \Gn{\Subst{!A}{t}{u}}.
\]
\end{prop}

\begin{proof}
तब हम फलन $\fn{hSubst}$ को आदिम पुनरावर्तन द्वारा निम्न प्रकार परिभाषित कर
सकते हैं:
\begin{multline*}
\begin{aligned}
\fn{hSubst}(x, y, z, 0) & = \emptyseq \\
\fn{hSubst}(x, y, z, i+1) & =
\end{aligned}\\
\begin{cases}
\fn{hSubst}(x, y, z, i) \concat y & \text{यदि $\fn{FreeOcc}(x, z, i)$} \\
\fn{append}(\fn{hSubst}(x, y, z, i), (x)_{i}) & \text{अन्यथा।}
\end{cases}
\end{multline*}
अब $\fn{Subst}(x, y, z)$ को $\fn{hSubst}(x, y, z, \len{x})$ के रूप में
परिभाषित किया जा सकता है।
\end{proof}

\begin{prop}
\ollabel{prop:free-for}
संबंध $\fn{FreeFor}(x, y, z)$, जो तभी लागू होता है जब ग्योडेल (G\"odel)
संख्या~$y$ वाला पद, ग्योडेल (G\"odel) संख्या~$x$ वाले सूत्र में ग्योडेल (G\"odel)
 संख्या~$z$ वाले चर के लिए !!{free for} हो, आदिम पुनरावर्ती है।
\end{prop}

\begin{proof} अभ्यास। \end{proof}

\begin{prob}
\olref[inc][art][sub]{prop:free-for} सिद्ध करें।
\end{prob}

\end{document}
